% ===== PRÉAMBULE CANONIQUE — à coller VERBATIM en tête de CHAQUE .tex =====
\documentclass[11pt,a4paper]{article}
\usepackage[utf8]{inputenc}\usepackage[T1]{fontenc}\usepackage{lmodern}
\usepackage[french]{babel}
\usepackage{amsmath,amssymb,mathtools}\usepackage{icomma}
\usepackage{siunitx}\sisetup{locale=FR}  % \qty{}{}, \si{}, \ang{}
\usepackage{xcolor}\usepackage{colortbl}
\usepackage{tikz}\usepackage{pgfplots}\pgfplotsset{compat=1.18}
\usepackage[siunitx,europeanresistors]{circuitikz} % circuits (élec.)
\usepackage[most]{tcolorbox}\usepackage{enumitem}\usepackage{array,booktabs}
\usepackage[a4paper,margin=2.2cm]{geometry}
\usetikzlibrary{arrows.meta,calc,angles,quotes,positioning,patterns,%
  backgrounds,shapes.geometric,decorations.pathmorphing,decorations.markings,intersections}
\definecolor{primary}{RGB}{222,49,99}\definecolor{primarydark}{RGB}{150,22,55}
\definecolor{defcol}{RGB}{0,114,178}\definecolor{methcol}{RGB}{52,92,148}
\definecolor{excol}{RGB}{0,135,90}\definecolor{attncol}{RGB}{200,40,55}
\tcbset{boxbase/.style={enhanced,breakable,boxrule=0.9pt,arc=2.5pt,
  left=3mm,right=3mm,top=2.3mm,bottom=2.3mm,fonttitle=\bfseries,coltitle=white}}
% --- boîtes TUTORIEL (noms EXACTS, ne pas renommer) ---
\newtcolorbox{cle}[1][]{boxbase,colback=defcol!6,colframe=defcol,
  title={\textsc{À retenir}\ifx\relax#1\relax\relax\else\textmd{ : \itshape #1}\fi}}
\newtcolorbox{methode}[1][]{boxbase,colback=methcol!7,colframe=methcol,sharp corners,
  title={\textsc{Méthode}\ifx\relax#1\relax\relax\else\textmd{ --- \itshape #1}\fi}}
\newtcolorbox{exemplebox}[1][]{boxbase,colback=excol!5,colframe=excol,
  title={\textsc{Exemple}\ifx\relax#1\relax\relax\else\textmd{ : \itshape #1}\fi}}
\newtcolorbox{piege}[1][]{boxbase,colback=attncol!6,colframe=attncol,
  title={\textsc{Piège}\ifx\relax#1\relax\relax\else\textmd{ : \itshape #1}\fi}}
\newtcolorbox{remarque}[1][]{boxbase,colback=black!4,colframe=black!45,
  title={\textsc{Remarque}\ifx\relax#1\relax\relax\else\textmd{ : \itshape #1}\fi}}
% --- boîtes EXERCICE / CORRIGÉ (noms EXACTS, auto-comptées) ---
\newtcolorbox[auto counter]{exo}[1][]{boxbase,colback=primary!4,colframe=primary,
  title={\textsc{Exercice}~\thetcbcounter\ifx\relax#1\relax\relax\else\textmd{ --- \itshape #1}\fi}}
\newtcolorbox[auto counter]{corrige}[1][]{boxbase,colback=excol!4,colframe=excol,
  title={\textsc{Corrigé}~\thetcbcounter\ifx\relax#1\relax\relax\else\textmd{ --- \itshape #1}\fi}}
\newcommand{\vv}[1]{\vec{#1}}\newcommand{\ds}{\displaystyle}
\newcommand{\R}{\mathbb{R}}\newcommand{\N}{\mathbb{N}}\newcommand{\Z}{\mathbb{Z}}
% (un \usepackage{fancyhdr}+en-tête est possible ; il est ignoré côté web)
% =========================================================================
\begin{document}
% THEME_TITLE: Le centre de gravité
\begin{center}
{\LARGE\bfseries\color{primarydark}Thème 3 — Le centre de gravité}\\[2pt]
{\color{primary}\itshape Statique — Fiche 3}
\end{center}
\vspace{1mm}

Un corps réel est un assemblage de très nombreuses particules, chacune attirée par la Terre. En
composant toutes ces petites forces parallèles, on obtient \emph{une seule} force — le \textbf{poids}
$\vv G=m\vv g$ — appliquée en un point bien précis : le \textbf{centre de gravité} $G$.

\begin{cle}[définition et propriété fondamentale]
Le \textbf{centre de gravité} $G$ est le point d'application de la résultante de toutes les forces de
pesanteur. La \textbf{ligne d'action du poids passe toujours par $G$}, qui est un \emph{point fixe} du
corps, \textbf{quelle que soit son orientation}.
\end{cle}

\begin{cle}[centre de gravité de formes homogènes]
Pour un corps \textbf{homogène et symétrique}, $G$ est « au centre », sur le(s) axe(s) de symétrie :
\begin{itemize}[leftmargin=1.5em,topsep=1pt,itemsep=1pt]
  \item \textbf{plaque triangulaire} : intersection des \emph{médianes} ;
  \item \textbf{cube / parallélépipède} : son \emph{centre} ;
  \item \textbf{cylindre} : le \emph{milieu} du segment joignant les centres des deux bases.
\end{itemize}
\end{cle}

\begin{methode}[$G$ d'un ensemble de masses alignées (de proche en proche)]
Les poids sont des forces parallèles : on les compose deux à deux. Pour deux masses $m_1$ (en $x_1$)
et $m_2$ (en $x_2$), $G$ vérifie $m_1\times\overline{G_1 G} = m_2\times\overline{G G_2}$, soit
\[
\boxed{\,x_G = \frac{m_1 x_1 + m_2 x_2}{m_1 + m_2}\,}
\]
($G$ plus proche de la masse la plus lourde). Pour trois masses ou plus, on compose $m_1$ et $m_2$,
puis le résultat avec $m_3$, et ainsi de suite.
\end{methode}

\begin{methode}[détermination expérimentale — méthode des suspensions]
Pour une plaque mince : (1) suspendre en $P_1$, tracer la \textbf{verticale} passant par $P_1$ ;
(2) suspendre en un autre point $P_2$, tracer la nouvelle verticale ; (3) $G$ est à l'\textbf{intersection}
des deux verticales. (Un objet suspendu pivote jusqu'à ce que $G$ soit à la verticale du point
d'accrochage.)
\end{methode}

\begin{piege}[$G$ peut se trouver hors de la matière]
Pour un anneau, un fer à cheval ou un boomerang, $G$ se situe dans le \emph{vide} entouré par la
matière. C'est pourtant bien là que s'applique, fictivement, tout le poids du corps.
\end{piege}

\begin{cle}[condition de stabilité]
Un corps posé reste en équilibre tant que la \textbf{verticale passant par $G$} tombe \emph{à
l'intérieur de sa base d'appui} (base de sustentation). Si cette verticale sort de la base, le corps
\textbf{bascule}. Plus $G$ est bas et la base large, plus l'équilibre est stable.
\end{cle}

\end{document}
